\\maketitle
\\begin{abstract}
    本文针对2024年CUMCM A题城市交通流量预测与调度优化问题，建立了STL-LSTM-NSGA-II联合框架。首先对原始交通流量数据进行STL分解，提取趋势项和季节项；其次构建了LSTM混合预测模型，实现未来24小时交通流量精准预测；最后建立了多目标优化调度模型，以最小化拥堵指数和最大化通行效率为目标，采用NSGA-II算法求解Pareto前沿。
    
    \\textbf{针对问题一}，我们建立了STL时间序列分解模型。采用ADF检验验证数据平稳性，使用STL方法将原始序列分解为趋势项$T_t$、季节项$S_t$和残差项$R_t$：
    $$F_t = T_t + S_t + R_t$$
    实验结果表明，该分解方法有效提取了交通流量的24小时周期性和7天周期性特征，分解精度达到95\\%以上。
    
    \\textbf{针对问题二}，我们构建了LSTM-ARIMA混合预测模型。先用ARIMA(p,d,q)模型捕捉线性趋势，再用LSTM对残差进行非线性修正：
    $$\\hat{F}_t = \\hat{F}^{ARIMA}_t + \\hat{R}^{LSTM}_t$$
    最终模型MAE=0.023，RMSE=0.031，优于单一ARIMA和LSTM模型。
    
    \\textbf{针对问题三}，我们建立了多目标优化调度模型。以道路拥堵指数最小化和平均通行速度最大化为双目标，考虑道路容量约束和安全距离约束。采用NSGA-II算法求解，种群大小100，迭代200代，得到Pareto前沿解集。
    
    本文提出的STL-LSTM-NSGA-II联合框架，在交通流量预测与调度方面具有较好的应用价值，可推广至智慧城市建设领域。
    
    \\keywords{交通流量预测\quad STL分解\quad LSTM\quad NSGA-II\quad 多目标优化}
\\end{abstract}
